\section{Method}\label{sec:method}
\subsection{Problem formulation}
Let $\mathcal{D}=\{(x_i,y_i)\}_{i=1}^{N}$ denote the training set and let
$f_\theta$ denote a model with parameters $\theta$. Define the input, output,
and constraints before introducing individual components.
\begin{equation}
  \theta^* = \argmin_{\theta}\frac{1}{N}\sum_{i=1}^{N}
  \ell\bigl(f_\theta(x_i),y_i\bigr)
  +\lambda\lVert\theta\rVert_2^2.
  \label{eq:objective}
\end{equation}
Equation~\eqref{eq:objective} is a notation example rather than a prescribed
training objective. Define the regularization strength $\lambda$ when it is used.

\subsection{Model components}
\subsubsection{Candidate representation}
Define each intermediate quantity and its dimension. Multi-line equations
should align at meaningful operators:
\begin{align}
  h &= f_\theta(x), & z_j &= g_\phi(c_j),\label{eq:representations}\\
  s_j &= \frac{h^\top z_j}{\tau}, &
  p_j &= \frac{\exp(s_j)}{\sum_k\exp(s_k)}.\label{eq:probability}
\end{align}
Here $c_j$ is a candidate and $\tau>0$ is a temperature. Explain whether the
representations share parameters, how they are normalized, and how decisions
are decoded.

\paragraph{Special cases}
Use a cases environment when a rule has distinct branches:
\begin{equation}
  \hat y = \begin{cases}
    \argmax_j p_j, & \max_j p_j\geq\delta,\\
    \text{abstain}, & \text{otherwise}.
  \end{cases}
\end{equation}

\subsection{Algorithms and mathematical statements}
Use numbered algorithm captions and explicit inputs and outputs.
Algorithm~\ref{alg:evaluate} illustrates a reproducible evaluation procedure.
\begin{algorithm}[ht]
\caption{Evaluation procedure}\label{alg:evaluate}
\begin{algorithmic}[1]
\REQUIRE Fixed model $f_\theta$, evaluation set $\mathcal{D}_{\mathrm{eval}}$
\ENSURE Aggregated metrics and per-example records
\STATE Initialize records $\mathcal{R}\gets\varnothing$
\FOR{each $(x,y)\in\mathcal{D}_{\mathrm{eval}}$}
  \STATE Compute $\hat y\gets f_\theta(x)$
  \STATE Measure the metric and execution cost
  \STATE Append $(\hat y,y,\mathrm{cost})$ to $\mathcal{R}$
\ENDFOR
\STATE \textbf{return} Aggregate$(\mathcal{R})$
\end{algorithmic}
\end{algorithm}

\begin{proposition}[Normalization example]\label{prop:normalization}
For finite real-valued scores $s_j$, the probabilities in
Equation~\eqref{eq:probability} satisfy $p_j>0$ and $\sum_jp_j=1$.
\end{proposition}
\begin{proof}
Exponentials of finite real numbers are positive. Dividing each term by their
sum yields positive values whose sum is one.
\end{proof}
Only retain mathematical statements that serve the paper's argument.
